\documentclass[../../../include/open-logic-section]{subfiles}

\begin{document}

\olfileid{sth}{ord-arithmetic}{add}
\olsection{الجمع الترتيبي}

لنفرض أننا نريد جمع~$\alpha$ و$\beta$. نضع نسخة من~$\beta$ عقب
نسخة من~$\alpha$ مباشرة. وإنما نأخذ \emph{نسختين} لأن
\olref[ordinals][basic]{ordinalsaresubsets} يبين أن إما
$\alpha \subseteq \beta$ أو~$\beta \subseteq \alpha$. ومعنى هذا في
الحدس أن نجتاز خطوات نمطها~$\alpha$، ثم نجتاز خطوات نمطها~$\beta$.
والترتيب الناتج جيد، فيتماثل، بموجب
\olref[ordinals][ordtype]{thmOrdinalRepresentation}، مع عدد ترتيبي
واحد. وذلك العدد هو الذي نسميه \emph{مجموع}~$\alpha$ و$\beta$.

هذه جملة الفكرة وراء الجمع الترتيبي. ولضبطها نبدأ بمعنى أخذ
\emph{نسخ} من المجموعات: نستعمل وسمين اعتباطيين، هما~$0$ و$1$، لنعلم
من أين جاء كل شيء.

\begin{defn}\ollabel{defdissum}
\emph{الجمع المنفصل} لـ$A$ و$B$ هو
$A \disjointsum B = (A\times \{0\}) \cup (B \times \{1\})$.
\end{defn}

ثم نضع ترتيبًا على أزواج الأعداد الترتيبية.

\begin{defn}
لأي أعداد ترتيبية~$\alpha_1, \alpha_2, \beta_1, \beta_2$، نقول إن:
\begin{align*}
	\tuple{\alpha_1, \alpha_2} \rlexless \tuple{\beta_1, \beta_2}\text{ إذا وفقط إذا كان }&
	\text{إما $\alpha_2 \in \beta_2$،}\\
	& \text{أو كان معًا $\alpha_2 = \beta_2$ و$\alpha_1 \in \beta_1$}
\end{align*}
\end{defn}

وهذا ترتيب \emph{معجمي معكوس}؛ إذ يبدأ بالعنصر الثاني، ثم يرجع إلى
الأول. وقد أردنا أن يكون~$\alpha \ordplus \beta$ نمط ترتيب نسخة من
$\alpha$ تأتي بعدها نسخة من~$\beta$، فنقول:

\begin{defn}\ollabel{defordplus}
لأي عددين ترتيبيين~$\alpha$ و$\beta$، يكون مجموعهما
$\alpha \ordplus \beta =
\ordtype{\alpha \disjointsum \beta, \rlexless}$.
\end{defn}
\noindent
وفي هذا تسامح يسير في الرمز؛ فحقنا أن نكتب
«$\Setabs{\tuple{x,y}\in(\alpha\disjointsum\beta)^2}{x \rlexless y}$»
مكان «$\rlexless$». وسنبقي هذا التسامح طلبًا للاختصار.

وتدل النتيجة الآتية، مع
\olref[ordinals][ordtype]{thmOrdinalRepresentation}، على سلامة
التعريف:

\begin{lem}\ollabel{ordsumlessiswo}
$\tuple{\alpha \disjointsum \beta, \rlexless}$ ترتيب جيد، لأي عددين
ترتيبيين~$\alpha$ و$\beta$.
\end{lem}

\begin{proof}
اتصال~$\rlexless$ على~$\alpha \disjointsum \beta$ بيّن. ولإثبات حسن
التأسيس، ثبت مجموعة غير خالية~$X \subseteq \alpha \disjointsum \beta$.
ولتكن~$Y$ مجموعة عناصر~$X$ التي إحداثيها الثاني أصغر ما يمكن، أي
$Y = \Setabs{\tuple{\gamma, i} \in X}{(\forall \tuple{\delta, j} \in X)i \leq j}$.
ثم خذ من~$Y$ العنصر الذي إحداثيه الأول أصغر.
\end{proof}
\noindent
فقد حصلنا على تعريف صريح حسن للجمع الترتيبي. ومن خواصه القريبة، مع
تذكر~$1 = \{0\}$ من~\olref[z][infinity-again]{defnomega}، ما يأتي:
\begin{prop}
$\alpha \ordplus 1 = \ordsucc{\alpha}$، لأي عدد ترتيبي~$\alpha$.
\end{prop}

\begin{proof}
اعتبر التماثل~$f$ من
$\ordsucc{\alpha} = \alpha \cup \{\alpha\}$ إلى
$\alpha\disjointsum1 = (\alpha \times \{0\})
\cup (\{0\} \times \{1\})$، المعطى بـ
$f(\gamma) = \tuple{\gamma, 0}$ لكل~$\gamma \in \alpha$، وبـ
$f(\alpha) = \tuple{0, 1}$.
\end{proof}
\noindent
وللجمع كذلك أوصاف عودية يسهل إثباتها.

\begin{lem}\ollabel{ordadditionrecursion}
لأي عددين ترتيبيين~$\alpha, \beta$، لدينا:
\begin{align*}
	\alpha\ordplus 0 &= \alpha\\
	\alpha \ordplus  (\beta\ordplus 1) &= (\alpha \ordplus  \beta) \ordplus  1\\
	\alpha  \ordplus  \beta &= \supstrict_{\delta < \beta}(\alpha \ordplus  \delta) && \text{إذا كان $\beta$ عددًا ترتيبيًا حديًا}
\end{align*}
\end{lem}

\begin{proof}
ننظر في الحالات على ترتيبها. أولًا:
\begin{align*}
	\alpha \ordplus  0
	&= \ordtype{\alpha \times \{0\}, \rlexless}\\
	&= \alpha\\
	\alpha \ordplus (\beta \ordplus  1)
	%&= \ordtype{(\alpha\times \{0\}) \cup (\ordtype{\beta \ordplus  1}\times \{1\}), \rlexless} \\
	&= \ordtype{(\alpha\times \{0\}) \cup (\ordsucc{\beta}\times \{1\}), \rlexless} \\
	&= \ordtype{(\alpha\times \{0\}) \cup (\beta \times \{1\}), \rlexless} \ordplus 1\\
	&= (\alpha \ordplus  \beta) \ordplus  1
\end{align*}
والآن لتكن~$\beta \neq \emptyset$ حدية. متى كان~$\delta < \beta$،
كان~$\delta\ordplus 1 < \beta$ أيضًا، ومن ثم كان
$\alpha \ordplus \delta$ مقطعًا ابتدائيًا حقيقيًا من
$\alpha \ordplus \beta$. فيكون~$\alpha \ordplus \beta$ حدًا أعلى
صارمًا للمجموعة
$X = \Setabs{\alpha \ordplus \delta}{\delta < \beta}$. ثم إذا كان
$\alpha \leq \gamma < \alpha \ordplus \beta$، كان بالضرورة
$\gamma = \alpha \ordplus \delta$ لبعض~$\delta < \beta$. ولهذا
$\alpha \ordplus \beta =
\supstrict_{\delta< \beta}(\alpha\ordplus \delta)$.
\end{proof}

غير أن في الأمر وجهًا آخر. فقد كان في وسعنا أن نعرّف الجمع الترتيبي
بـ\emph{مجرد} مبرهنة العودية المتناهية التجاوز، فنقرر معادلات
العودية الواردة في~\olref{ordadditionrecursion}، مع استعمال
«$\ordsucc{\beta}$» مكان «$\beta \ordplus 1$».

فللعمليات على الأعداد الترتيبية طريقان: طريق \emph{تركيبي} نبني فيه
مجموعة مرتبة ترتيبًا جيدًا ثم نأخذ نمط ترتيبها، وطريق \emph{عودي}
نقرر فيه معادلات العودية. وإذا أحكم العمل اتفق الطريقان؛ لأن مبرهنة
العودية المتناهية التجاوز تضمن أن تلك المعادلات تعين أعدادًا ترتيبية
\emph{فريدة}.

والحساب الترتيبي يشبه، من وجوه كثيرة، حساب الأعداد الطبيعية؛ ومن ذلك:

\begin{lem}\ollabel{ordinaladditionisnice}
إذا كانت~$\alpha, \beta, \gamma$ أعدادًا ترتيبية، فإن:
\begin{enumerate}
	\item\ollabel{ordaddition1} إذا كان~$\beta < \gamma$، فإن
	$\alpha \ordplus \beta < \alpha \ordplus \gamma$؛
	\item\ollabel{ordaddition2} إذا كان~$\alpha \ordplus \beta =
	\alpha\ordplus \gamma$، فإن~$\beta = \gamma$؛
	\item\ollabel{ordaddition3} $\alpha \ordplus (\beta \ordplus
	\gamma) = (\alpha \ordplus \beta) \ordplus \gamma$، أي إن الجمع
	تجميعي؛
	\item\ollabel{ordaddition4} إذا كان~$\alpha \leq \beta$، فإن
	$\alpha \ordplus \gamma \leq \beta \ordplus \gamma$.
\end{enumerate}
\end{lem}

\begin{proof}
نبرهن~\olref{ordaddition3} ونكل الباقي إلى التمرين. نستعمل الاستقراء
البسيط المتناهي التجاوز على~$\gamma$، مع
\olref{ordadditionrecursion}. فإذا~$\gamma = 0$:
\[
(\alpha \ordplus  \beta) \ordplus  0 = \alpha \ordplus  \beta  = \alpha \ordplus  (\beta \ordplus  0)
\]
وإذا~$\gamma = \delta\ordplus 1$، فافترض استقرائيًا أن
$(\alpha \ordplus  \beta) \ordplus  \delta =
\alpha \ordplus  (\beta \ordplus \delta)$؛ ثم استعمل
\olref{ordadditionrecursion} ثلاث مرات:
\begin{align*}
	(\alpha \ordplus  \beta) \ordplus  (\delta \ordplus  1) & = ((\alpha \ordplus  \beta) \ordplus  \delta)\ordplus 1\\
	& = (\alpha \ordplus  (\beta \ordplus  \delta)) \ordplus  1\\
	& = \alpha \ordplus  ((\beta \ordplus  \delta)\ordplus 1)\\
	& = \alpha \ordplus  (\beta \ordplus  (\delta\ordplus 1))
\end{align*}	
وأما إذا كانت~$\gamma$ عددًا ترتيبيًا حديًا، فافترض استقرائيًا أنه
متى كان~$\delta \in \gamma$، فإن
$(\alpha \ordplus  \beta) \ordplus  \delta =
\alpha \ordplus  (\beta \ordplus  \delta)$؛ وعندئذ:
\begin{align*}
	(\alpha \ordplus  \beta) \ordplus  \gamma & = \supstrict_{\delta < \gamma}((\alpha \ordplus  \beta) \ordplus  \delta) \\
	&= \supstrict_{\delta < \gamma}(\alpha \ordplus  (\beta \ordplus  \delta))\\
	&= \alpha \ordplus  \supstrict_{\delta < \gamma}(\beta \ordplus  \delta)\\
	& = \alpha \ordplus  (\beta \ordplus  \gamma)
\end{align*}
\end{proof}

\begin{prob}
أثبت بقية~\olref[sth][ord-arithmetic][add]{ordinaladditionisnice}.
\end{prob}

بهذه الخواص يبدو الجمع الترتيبي مألوفًا، غير أنه يفارق جمع الأعداد
الطبيعية في وجه جوهري.

\begin{prop}\ollabel{ordsumnotcommute}
الجمع الترتيبي {ليس} تبادليًا؛
$1 \ordplus \omega = \omega < \omega \ordplus 1$.
\end{prop}

\begin{proof}
لاحظ أن~$1 \ordplus \omega = \supstrict_{n < \omega} (1 \ordplus n)
= \omega \in \omega \cup \{\omega\} = \ordsucc{\omega} = \omega
\ordplus 1$.
\end{proof}
\noindent
وقد يفجأ هذا أول النظر، ولكنه لا ينبغي أن يفجأ. فإن
$1 \ordplus \omega$ هو نمط الترتيب الحاصل بإدخال عنصر زائد
\emph{قبل} الأعداد الطبيعية كلها. وكما في حجة فندق هيلبرت في
\olref[sfr][infinite][hilbert]{sec}، لا يغير ذلك العنصر الأول نمط
الترتيب الكلي في الحدس. وأما~$\omega \ordplus 1$ فهو نمط الترتيب
الحاصل بإلحاق عنصر \emph{بعد} الأعداد الطبيعية كلها؛ وبين الحالين
فرق بين.

\end{document}
